\documentclass[instructions.tex]{subfiles}
\begin{document}
 
\chapter{De la mort. Il faut s’y préparer.}
 
\primo Il n’est rien de plus certain et de plus rigoureux que la mort; il n’est rien de plus incertain que ses circonstances, et rien de plus immuable que la destinée du mourant; il faut donc s’y préparer.
 
Il n’est rien de plus certain que la mort, c’est un arrêt porté par le Tout-Puissant\footnote{Statutum est hominibus semel mori. Hebr.. 27. 9.— Mort à non putatis. Luc. 12. 40.}. Il n’est rien de plus rigoureux ici-bas; mourir, c’est quitter la terre pour toujours; c’est entrer dans l’éternité, sans savoir où l’on va; c’est perdre la vie avec tout ce qu’on possede. La mort est un adieu général à tout ce qu’on a de plus cher : parens, amis, richesses, honneurs, emplois, sans exception et sans retour. A la mort tout nous quitte; il n’y a que le vice ou la vertu qui nous accompagne.
 
Si la pensée de la mort est si insupportable à ceux qui sont attachés à la vie, aux voluptueux et aux riches, combien terrible et rigoureux en doit être le coup ! O moment effroyable ! qui leur découvre enfin la vanité des biens et des plaisirs qu’ils vont perdre pour toujours.
 
S’il n’est rien de plus certain et de plus rigoureux que la mort, il n’est rien de plus incertain que le temps et les circonstances. Mourrez-vous dans la grâce ou dans le péché ? Mourrez-vous d’une maladie lente ou subitement ? Aurez-vous le temps de vous préparer, ou ne serez-vous point surpris dans peu de jours, cette nuit ou dans un moment ? Personne ne peut le dire; vous ne pouvez être assuré d’un seul instant. Ce qui est certain, c’est que vous mourrez, et que vous mourrez quand vous y penserez le moins\footnote{Qua hora non putatis. lux. 12.40}.
 
Mais ce qui est bien plus sérieux, c’est que la destinée du mourant est immuable. La mort décide de tout pour l’éternité. Si je meurs dans le péché mortel, je serai réprouvé et malheureux pour toujours. Si je meurs dans la grâce de Dieu, je serai sauvé et heureux pour jamais. Où l’arbre tombera, il restera\footnote{In quocumque loco ceciderit, ibi erit. Eccl. 11,}.
 
\secundo Puisque la mort est inévitable, et que ce moment terrible doit décider de votre salut, vous devez donc vous y préparer, et profiter du temps qui vous reste. Si un prisonnier qui, dans quelques jours, doit être condamné au gibet, s’échappait de prison, s’arrêterait-il dans sa course ? s’amuserait-il à des bagatelles ? ne s’observerait-il pas, crainte d’être surpris dans sa fuite ? Mais si l’on permettait à ce malheureux fugitif de prendre dans un trésor, pendant un jour, tout ce qu’il voudrait pour obtenir sa grâce, perdrait-il un moment d’un jour si précieux ? Et vous, misérables mortels, qui avez si souvent échappé à l’enfer ! vous que dans peu la mort citera devant votre juge, à quoi vous amusez-vous ? Vous pouvez, pendant le peu de jours que vous voyagez sur la terre, satisfaire à Dieu et vous enrichir pour le ciel\footnote{Parare possunt æternam vitam dies pauci. Euch.}. Mais à quoi, hélas ! employez-vous des jours si précieux et si courts ?
 
Faites-vous réflexion que vous ne savez ni le terme que Dieu a fixé à votre vie, ni le nombre des grâces qu’il vous destine, ni le nombre des péchés qu’il veut souffrir de vous ? Pourquoi donc négligez-vous de profiter du temps ? Ne vous fiez pas à votre santé, les apparences d’une longue vie sont trompeuses : la mort peut vous enlever en un moment; votre temps s’écoule. Hélas ! combien en avez-vous déjà perdu ! Employez du moins saintement le peu qui vous reste; le ciel ou l’enfer doit être la récompense ou le châtiment de l’usage que vous en ferez. Un officier de l’empereur Charles-Quint, pénétré de cette importante vérité, lui demanda la permission de se retirer, alléguant pour raison, qu’entre les affaires du monde et le jour de la mort il doit y avoir du temps. L’empereur, touché de cet exemple, céda l’empire à son frère, ses royaumes à son fils, et se retira du monde pour se disposer à la mort. Peut-on prendre trop de précautions pour se préparer à un passage qu’on ne fait qu’une fois, et qui, dans un moment, décide de tout et pour toujours ?
 
\section{Résolutions}
 
\primo Il faut que je me détache dès à présent de tout ce que je dois quitter à la mort. 

\secundo Et puisque je ne puis éviter ce terrible passage, qui décidera infailliblement de mon éternité, il faut que je commence dès ce moment à m’y préparer, en vivant tous les jours comme si chacun d’eux devait être le dernier de mes jours. 
\prayer{O mon Dieu ! faites-moi la grâce de mourir dans votre saint amour.}
 
\end{document}
